\documentclass[11pt]{article}
\usepackage[margin=1in]{geometry}
\usepackage{amsmath,amssymb,amsthm}
\newtheorem{theorem}{Theorem}
\newtheorem{lemma}{Lemma}
\newtheorem{corollary}{Corollary}
\title{Vertical Core Slices Remove the Interior-Threshold Hypothesis}
\author{}\date{October 5, 2026 --- paper-proof draft}
\begin{document}\maketitle

Fix $0<w<\pi/2$, $0<a<w/2$, $b=w-a$. Use the support functions $p,k$, corner $\gamma=(p-1)u_t+(k-1)v_t$, fan $F_w$, and functional $Q$ from \texttt{lifted-contact-certificate.tex}. Put $A=p-1$, $B=k-1$ and
\[
 H_R=\{z:z\cdot u_a\ge A(a)\},\qquad
 H_L=\{z:z\cdot v_b\ge B(b)\}.
\]
The earlier lifting proof assumes every interior threshold is positive to justify radial contractions. This note replaces that step by vertical slices. No positivity of $A(t),B(t)$ away from the two selected cuts is required.

\begin{theorem}[Lifting with exactly the available geometry]\label{thm:lift}
Suppose $K$ is a normalized cap, its corner is continuously differentiable, $f=k-p'>1$, $g=p+k'>1$ on $(0,w)$, and $\gamma([0,w])\subset F_w$. Suppose also
\[
 A(a)>0,\quad B(b)>0,\qquad |N_w(K)\cap H_R\cap H_L|=0.
\]
Then there are admissible lifted functions $\beta,\eta$ in the obstacle domain of \texttt{lifted-contact-certificate.tex} such that
\[
 \mathcal A_w(K)\le Q(p,k,\beta,\eta).
\]
\end{theorem}

\begin{proof}
The arm formula gives
\[
 \gamma_x'=-(f-1)\cos t-(g-1)\sin t<0.
\]
The right cut meets the floor at $R=(A(a)/\cos a,0)$; the left cut meets the other fan ray at $L=B(b)v_w/\cos a$. These points have respectively positive and negative abscissae. Furthermore
\[
 R_x-\gamma_x(a)=\gamma_y(a)\tan a\ge0,
\]
and $\gamma(b)-L$ is a nonnegative multiple of $u_b$, because its scalar product with $u_w$ is nonnegative. Thus
\[
 R_x\ge\gamma_x(a)>\gamma_x(b)\ge L_x.
\]
The piecewise curve from $R$ to $\gamma(a)$ along the cut, then along $\gamma$ to $\gamma(b)$, then along the other cut to $L$, is a continuous graph $y=Y(x)$ over $[L_x,R_x]$. Degenerate end segments can be omitted. It lies in the fan: the end segments join fan points and the middle curve is in the fan by hypothesis. The fan's lower graph is
\[
 \ell(x)=\max\{0,-x\cot w\}.
\]
Let $D_{a,b}$ be the region between $Y$ and $\ell$.

We verify its inclusion in the niche by three kinds of vertical slice. For $\gamma_x(b)<x<\gamma_x(a)$ choose the unique $t\in(a,b)$ with $\gamma_x(t)=x$. Every point $z=(x,y)$ with $\ell(x)\le y<\gamma_y(t)$ is obtained by moving vertically down from $\gamma(t)$. Both wall coordinates decrease strictly, since $\sin t,\cos t>0$, so $z\in Q^-_K(t)\cap F_w$.

For $\gamma_x(a)\le x<R_x$, the right cut graph has height
$Y(x)=\gamma_y(a)-(x-\gamma_x(a))\cot a\le\gamma_y(a)$.
A point strictly below this graph satisfies $z\cdot u_a<A(a)$ and
\[
 (z-\gamma(a))\cdot v_a
 =-(x-\gamma_x(a))\sin a+(y-\gamma_y(a))\cos a<0.
\]
It therefore belongs to the niche at angle $a$, even if $B(a)$ is negative. On the left end segment $x\le\gamma_x(b)$ and
$Y(x)=\gamma_y(b)+(x-\gamma_x(b))\tan b\le\gamma_y(b)$.
Points below it have $z\cdot v_b<B(b)$ and $(z-\gamma(b))\cdot u_b<0$, so belong to the niche at angle $b$. Thus $D_{a,b}\subset N_w(K)$ up to its graph and fan boundary, all of which have area zero.

This core lies below both cuts. On the middle curve $\gamma(t)\cdot u_a$ decreases for $t\ge a$, while $\gamma(t)\cdot v_b$ increases for $t\le b$, by the arm formula. The right floor point satisfies $R\cdot v_b\le0<B(b)$, and the left fan point satisfies $L\cdot u_a<0<A(a)$. Hence both end segments satisfy the other cut inequality as well. Lowering the ordinate decreases both cut projections. The entire core consequently avoids the interiors of $H_R,H_L$.

For completeness the exterior construction requires no removed hypothesis. Set
\[
 W_R=\{y\ge0,z\cdot u_a\ge A(a)\},\quad
 E_R=W_R\cap\bigcap_{a\le t\le w}\{z:z\cdot u_t\ge p(t)-1\}.
\]
Its infimum support $\beta(t)=\inf_{z\in E_R}z\cdot u_t$ on $[a,\pi/2]$ is finite and Lipschitz, satisfies $\beta(a)=A(a)$, $\beta(\pi/2)=0$, and majorizes $p(t)-1$ for $a\le t\le w$. Indeed the floor intersects $E_R$ at a finite abscissa, and its cut line intersects $E_R$ at a finite height: on that line the required later heights are bounded using the Lipschitz bound on $p$ and $\sin(t-a)/(t-a)\to1$. The lower boundary between these intersections is compact and supplies all these infimum supports.

The wedge lies in $F_w$. If $z\in W_R\setminus E_R$ violates the inequality at $s>a$, write $z-\gamma(s)=Xu_s+Yv_s$. Then $X<0$. Since $\gamma(s)\cdot u_a\le A(a)\le z\cdot u_a$, one has $X\cos(s-a)-Y\sin(s-a)\ge0$, whence $Y<0$. Thus $W_R\setminus E_R\subset N_w(K)\cap H_R$. Reflection gives $\eta$ and the analogous left region. The cut-separation assumption makes these two regions disjoint up to null sets, and they are disjoint from the core.

The Green area formula for the core graph and its two straight end segments is
\[
 |D_{a,b}|=I(\gamma|_{[a,b]})+\frac12[A(a)B(a)+A(b)B(b)]
 +\frac12[A(a)^2+B(b)^2]\tan a.
\]
This is a graph area identity; the mixed products need not individually be positive. The exterior area formula is
\[
 |W_R\setminus E_R|=-\frac12\int_a^{\pi/2}(\beta^2-\beta'^2)-\frac12A(a)^2\tan a,
\]
with its reflected analogue. This is the same bounded convex-boundary support-area identity as in the original lift. Adding the three disjoint areas cancels the tangent-square terms and gives the niche lower bound defining $Q$. Subtract it from the cap support-area formula to obtain the theorem.
\end{proof}

\begin{corollary}[The lift covers every larger-angle maximizer]
For every $1\le w<\pi/2$, every positive global maximizing cap and every $0<a\le1/12$ admit the bound in Theorem~\ref{thm:lift}.
\end{corollary}
\begin{proof}
The fixed-angle curvature and arm certificates give the corner regularity and strict arm inequalities. \texttt{global-cut-geometry.tex} gives fan containment, the cut thresholds at least $13/192$, and empty intersection of the niche with the two cut half-planes. These are exactly the hypotheses above.
\end{proof}

\paragraph{Scope.}
This closes a hypothesis mismatch in the all-angle use of the lift. It does not prove that the relaxed lifted maximum is attained by a feasible sofa, nor that the inequality is sharp for an arbitrary cut. Those are separate contact/exposure obligations. In particular, cut separation alone is not a proof of an all-angle optimizer theorem. No CI, Lean compilation, or TeX compilation was used.
\end{document}
